\documentclass[]{fim-uhk-thesis}
%\documentclass[english]{fim-uhk-thesis} % without assignment - for the work start to avoid compilation problem
% * Je-li práce psaná v anglickém jazyce, je zapotřebí u třídy použít
%   parametr english následovně:
%   If thesis is written in English, it is necessary to use
%   parameter english as follows:
%      \documentclass[english]{fim-uhk-thesis}
% * Je-li práce psaná ve slovenském jazyce, je zapotřebí u třídy použít
%   parametr slovak následovně:
%   If the work is written in the Slovak language, it is necessary
%   to use parameter slovak as follows:
%      \documentclass[slovak]{fim-uhk-thesis}
% * Je-li práce psaná v anglickém jazyce se slovenským abstraktem apod.,
%   je zapotřebí u třídy použít parametry english a enslovak následovně:
%   If the work is written in English with the Slovak abstract, etc.,
%   it is necessary to use parameters english and enslovak as follows:
%      \documentclass[english,enslovak]{fim-uhk-thesis}

% Základní balíčky jsou dole v souboru šablony fim-uhk-thesis.cls
% Basic packages are at the bottom of template file fim-uhk-thesis.cls
% zde můžeme vložit vlastní balíčky / you can place own packages here



%---rm---------------
\renewcommand{\rmdefault}{lmr}%zavede Latin Modern Roman jako rm / set Latin Modern Roman as rm
%---sf---------------
\renewcommand{\sfdefault}{qhv}%zavede TeX Gyre Heros jako sf
%---tt------------
\renewcommand{\ttdefault}{lmtt}% zavede Latin Modern tt jako tt

% vypne funkci šablony, která automaticky nahrazuje uvozovky,
% aby nebyly prováděny nevhodné náhrady v popisech API apod.
% disables function of the template which replaces quotation marks
% to avoid unnecessary replacements in the API descriptions etc.
\csdoublequotesoff



\def\UrlBreaks{\do\/\do-}
\usepackage{breakurl}


% =======================================================================
% balíček "hyperref" vytváří klikací odkazy v pdf, pokud tedy použijeme pdflatex
% problém je, že balíček hyperref musí být uveden jako poslední, takže nemůže
% být v šabloně
% "hyperref" package create clickable links in pdf if you are using pdflatex.
% Problem is that this package have to be introduced as the last one so it
% can not be placed in the template file.

  \usepackage[unicode,colorlinks,hyperindex,plainpages=false,urlcolor=black,linkcolor=black,citecolor=black]{hyperref}
  \definecolor{links}{rgb}{0,0,0}
  \definecolor{anchors}{rgb}{0,0,0}
  \def\AnchorColor{anchors}
  \def\LinkColor{links}
  \def\pdfBorderAttrs{/Border [0 0 0] } % bez okrajů kolem odkazů / without margins around links
  \pdfcompresslevel=9

% Řešení problému, kdy klikací odkazy na obrázky vedou za obrázek
% This solves the problems with links which leads after the picture
\usepackage[all]{hypcap}

\newcommand{\source}[1]{\caption*{Zdroj: {#1}} }

\projectinfo{
  %Prace / Thesis
  project={DP},            %typ práce BP/SP/DP/DR  / thesis type (SP = term project)
  % enter year of submission
  year={2023},             % rok odevzdání / year of submission
  date=\today,             % datum odevzdání / submission date
  %Nazev prace / thesis title
  title.cs={Návrh testovacího prostředí za využití nástrojů ATT\&CK},  % název práce v češtině či slovenštině (dle zadání) / thesis title in czech language (according to assignment)
  title.en={Design of test environment using ATT\&CK tools}, % název práce v angličtině / thesis title in english
  %Autor / Author
  author.name={Jiří},   % jméno autora / author name
  author.surname={Bönsch},   % příjmení autora / author surname
  author.title.p={Bc.},   %titul před jménem (nepovinné) / title before the name (optional)
  % enter field
  author.field={Aplikovaná Informatika},   % obor
  %Ustav / Department
  department={KIT}, % doplňte příslušnou zkratku dle ústavu na zadání: DEF / fill in appropriate abbreviation of the department according to assignment: DEF
  % Školitel / supervisor
  %  enter supervisor
  supervisor.name={Josef},   % jméno školitele / supervisor name
  supervisor.surname={Horálek},   % příjmení školitele / supervisor surname
  supervisor.title.p={Mgr.},   %titul před jménem (nepovinné) / title before the name (optional)
  supervisor.title.a={Ph.D.},
  % Anotace
  % enter czech annotation
  annotation.cs={
    Cílem práce bylo poukázat na současný stav kybernetické bezpečnosti a vytvořit jednoduché a srozumitelné testovací prostředí spolu s návodem na jeho použití.
    Byla provedena rešerše literatury, kde jsou zkoumány nejnutnější informace nutné k pochopení kybernetické bezpečnosti, hrozby v podobě druhů útočníků a nejčastější typů útoků.
    Také jsou představeny základní principy testování, které jsou poté vyzkoušeny pomocí vytvořeného testovacího prostředí na simulované infrastruktuře s využitím scénářů, jenž napodobují reálné chování útočníka.
    Aby měla práce relevantní dopad, je simulovaná infrastruktura založena na reálné infrastruktuře využívané ve skutečné organizaci.
    Výsledkem je zhodnocení scénářů, ukázka silných i slabých stránek infrastruktury a postupy, jak sestavit a upravit testovací prostředí pro potřeby testování v praxi.
  },
  %  enter english annotation
  annotation.en={
    The aim of this work was to highlight the current state of cyber-security and create a simple and understandable testing environment along with instructions for its use.
    A literature review was conducted, examining the essential information necessary for understanding cyber-security, threats posed by different types of attackers, and the most common types of attacks.
    Basic testing principles were also introduced and tested using the created testing environment on simulated infrastructure using scenarios that mimic real attacker behavior.
    To ensure the work has relevant impact, the simulated infrastructure is based on real infrastructure used in an actual organization.
    The results include an evaluation of the scenarios, demonstration of strengths and weaknesses of the infrastructure, as well as procedures for creating and modifying the testing environment for practical testing purposes.
},
  % WRITER: enter declaration if you want something else
  declaration={
    Prohlašuji, že jsem diplomovou práci zpracoval samostatně a s použitím uvedené literatury.
  },
  % enter acknowledgment, you can only change name of the supervisor
  acknowledgment={
    Tímto bych chtěl poděkovat vedoucímu diplomové práce Mgr. Josefu Horálkovi, Ph.D. za odborné vedení práce, poskytnuté rady a přínosné konzultace.
    Zároveň bych chtěl poděkovat rodině a přátelům za jejich pomoc a podporu.
  },
%(external submitter, consultant, etc.).},
  faculty={FIM},
  department={KIT}
}

% řeší první/poslední řádek odstavce na předchozí/následující stránce
% solves first/last row of the paragraph on the previous/next page
\clubpenalty=10000
\widowpenalty=10000

% @SEE http://ftp.cvut.cz/tex-archive/macros/latex/contrib/acro/acro-manual.pdf

\DeclareAcronym{IoT}{
  short=IoT,
  long=Internet of Things
}
\DeclareAcronym{DoS}{
  short=DoS,
  long=Denial-of-Service
}
\DeclareAcronym{DDoS}{
  short=DDoS,
  long=Distributed Denial-of-Service
}
\DeclareAcronym{RDoS}{
  short=RDoS,
  long=Ransomware Denial-of-Service
}
\DeclareAcronym{DNS}{
  short=DNS,
  long=Domain Name System
}
\DeclareAcronym{RaaS}{
  short=RaaS,
  long=Ransomware as a Service
}
\DeclareAcronym{MaaS}{
  short=MaaS,
  long=Malware as a Service
}
\DeclareAcronym{RTA}{
  short=RTA,
  long=Red Team Automation
}
\DeclareAcronym{ENISA}{
  short=ENISA,
  long=European Union Agency for Cybersecurity
}
\DeclareAcronym{OSINT}{
  short=OSINT,
  long=Open Source Inteligence
}
\DeclareAcronym{EU}{
  short=EU,
  long=Evropská unie
}
\DeclareAcronym{NIST}{
  short=NIST,
  long=National Institute of Standards and Technology
}
\DeclareAcronym{NSA}{
  short=NSA,
  long=National Security Agency
}
\DeclareAcronym{RDP}{
  short=RDP,
  long=Remote Desktop Protocol
}
\DeclareAcronym{MFA}{
  short=MFA,
  long=Multifactor Authentication
}
\DeclareAcronym{VLAN}{
  short=VLAN,
  long=Virtual Local Area Network
}
\DeclareAcronym{PhaaS}{
  short=PhaaS,
  long=Phishing as a Service
}
\DeclareAcronym{IKE}{
  short=IKE,
  long=Internet Key Exchange
}
\DeclareAcronym{USB}{
  short=USB,
  long=Universal Serial Bus
}

% začátek dokumentu
\begin{document}
  %fix if not enought pages -> \linespread{1.25}
  \linespread{1.25}
  % Vysazeni titulnich stran / Typesetting of the title pages
  % ----------------------------------------------
  \maketitle
  % Obsah
  % ----------------------------------------------
  \setlength{\parskip}{0pt}

  \thispagestyle{empty}
  {\hypersetup{hidelinks}\tableofcontents}

  % vynechani stranky v oboustrannem rezimu
  % Skip the page in the two-sided mode
  \iftwoside
    \cleardoublepage
  \fi

  % Text prace / Thesis text
  % ----------------------------------------------
  \setcounter{page}{0}
  \section{Úvod}

% TODO WRITER: text

Zde vysvětlit problémovou situaci a otázky, které se budou v bakalářské/diplomové práci řešit.

\section{Cíl práce}

% TODO WRITER: text

Smysl a účel, výzkumné otázky.

%teoretická část
\section{Metodika zpracování}

% TODO WRITER: text

Cíle, hypotézy/ výzkumné otázky, způsob hledání odpovědí na výzkumné otázky včetně metodiky vlastního výzkumu/šetření, literární rešerše.

\section{Úvod do problematiky Cyber-security}
%todo writer clean whole text, add as needed
V dnešní době díky rozvoji moderních technologií tráví velká část lidské populace značné množství času interakcemi s těmito technologiemi.
Ať už se jedná o Internet, \firstac{IoT}, mobilní bankovnictví nebo sociální sítě, lidé se na technologií stávají čím dál více závislí.
Technologie nám usnadňují všední život, dovolují práci z domova nebo přináší příležitosti, které jen pár let předem neexistovaly.
Tento rozvoj bohužel není jen kladný.
Rozvoj se projevil i hrozbami, kterým populace čelí.
Některé hrozby, jako malware jsou pro lidstvo nové.
Jiné jsou staré jako lidstvo samo a pouze se novým okolnostem přizpůsobily.
Jedním příkladem za všechny je Sociální inženýrství.

\paragraph{}
Hrozbám nečelí jen samotní lidé ale i firmy a dokonce státy.
Státy na celém světě se staly na kybernetickém světe zcela závislé.\cite{LI20218176}
Jakékoli jeho selhání může ohrozit jejich funkčnost a odvázání se od něj je již nemožné.
Existuje nezměrné množství hrozeb.
Je jasné, že v prostředí kde mezi sebou soupeří státy je značné úsilí věnováno jak zabezpečení tak útoku.
Rozdíly mezi cyber-útokem, cyber-zločinem a cyber-válkou jsou ve většině případů rozděleny pouze dle účastníků.\cite{LI20218176}
Způsob útoku či hrozby na zatřídění do předchozích kategorií klade minimální vliv, zohledněn bývá pouze výsledek.


\section{Relevance}
Relevantnost problematiky cyber-bezpečnosti potvrzuje nejen množství akademických textů které zkoumají tuto problematiku, ale také stále větší zájem všeobecné veřejnosti.
Hlavním důvodem je vzrůstající povědomí populace, způsobené medializací útoků sociálního inženýrství nebo ransomware které populaci v poslední době ovlivňovaly.
V době psaní je v popředí například ransomware útok na dětskou nemocnici SickKids.\cite{bleep_sickkids_ransom}
I když Ransomware skupina \textit{LockBit} poskytla nemocnici \textit{decryptor} zdarma a s omluvou, nemění to nic na vážnosti situace.
Je totiž důležité si uvědomit, že takovýto útok nastal.
Byl schopen omezit optimální funkčnost nemocnice zašifrováním důležitých interních systému.
A bohužel není jediným ransomware útokem zaměřeným na zdravotnictví.
Je tedy možné předpokládat, že nastanou i další útoky na nemocnice.
A příští útok nemusí být veden skupinou, která se řídí silným morálním kompasem.
Důležité je na závěr poznamenat, že tento silný morální kompas platí jen na omezení funkčnosti které může vést ke smrti.
Nevztahuje se na méně kritické instituce nebo ani na prodej již získaných zdravotních dat na černém trhu.

\paragraph{}
Je tedy důležité věnovat se i nadále cyber-bezpečnosti.
Jsou hrozby, které jsou dobře prostudované a víme, jak se jim efektivně bránit.
Jsou hrozby, kterým je sice obtížné se bránit, ale jsou omezené rozsahem a tak je obrana uskutečnitelné.
Ale také jsou hrozby, na které obrana zatím neexistuje.
Tyto hrozby je nutné najít a zkoumat, aby bylo možné ochranu vytvořit.

\subsection{Studie trendů Cyber-security}
Dlouhodobě můžeme v rámci Cyber-security rozeznávat několik klíčových zaměření.
Zajímavé jsou změny v prioritách zaměření v ohledu na čas, které můžeme zkoumat na literatuře.
Velké změny lze zaznamenat především před a po pandemii Covidu-19.\cite{KUMAR2022102821}
Jako i u jiných odvětví můžeme dělit celou sféru na 2 části, část akademickou a část veřejnou.
Je jasné, že zaměření těchto dvou částí budou lehce rozdílná, proto je nutné dívat se na ně odděleně i když se v mnoha ohledech překrývají a ovlivňují.
Důvod proč zkoumáme rozdíl před a po Covidu-19 je rapidní navýšení adopce digitálních technologií, které v této době proběhlo.
Ztížené pracovní podmínky vyžadovaly zavádění nových technologií.
To znamená nová bezpečnostní rizika vzniklá nejen z těchto technologií, ale i z přístupu jejich využívání.
Je značný rozdíl v bezpečnosti práce uživatele, nachází-li se v uzavřené síti uvnitř společnosti nebo při připojování na dálku při práci z domova.

\paragraph{Akademická literatura}
V akademické sféře je možné vidět zaměření výzkumu.
Dlouhodobě populárními tématy jsou jsou \textit{cyber risk management, detekce malware} nebo \textit{systémy detekce přístupu}.\cite{KUMAR2022102821}
Tato témata je možné nalézt jak před, tak i po pandemii Covidu-19.
Témata, které byla mnohem více zkoumána před pandemií jsou například využití \textit{Machine learningu}, a to jak pro detekci hrozeb tak hledání nových útočných vektorů.
Také \textit{Blockchain} a \textit{Cryptoměna} jsou témata, která byla před pandemií v oblasti zájmu akademiků mnohem více, nežli po pandemii.\cite{KUMAR2022102821}
Pandemie způsobila velký zvrat v myšlení.
Vedlo to ke zkoumání témat, které bylo nutné okamžitě aplikovat nebo která byla nutné pro budoucí fungování společnosti.
Objevují se zde mnohem více témata zaměřená na \textit{Zdravotnictví, Bankovní sektor} nebo \textit{zranitelnosti dodavatelského řetězce}.\cite{KUMAR2022102821}
Všechna tato témata mají jasný vznik v okolnostech pandemie a jejich zkoumání je prováděno i roky po pandemii samotné.

\paragraph{Všeobecná literatura}
Všeobecná literatura poukazuje na zájmy a zvědavost lidu, proto zde lze odvodit zaměření ale také problémy populace.
Před pandemií byly nejčastější témata zaměřená na \textit{finančně zaměřené kybernetické útoky}.
Ty spolu s běžnými tématy jako \textit{běžné bezpečnostní slabiny, úniky dat} a \textit{malware incidenty} vzbuzovaly největší zájem médií.\cite{KUMAR2022102821}
Pandemie rapidně změnila zájem lidu.
Zaměření se přesunulo na útoky, které se během pandemie značně rozšířily.
Jsou jimi hlavně \textit{Sociální inženýrství} nebo útoky \textit{ransomware}.
Také útoky na \textit{Zdravotnictví} byli mnohem častěji reportovány medii.
Tématem, které je v obou obdobích významné je \textit{malware}.
Jeho primární zaměření se však liší.
Před pandemií byla literatura zaměřena hlavně na ovlivnění správného chodu firem.
Po pandemii se psalo spíše o dopadech na digitální infrastrukturu jako jsou \textit{služby v cloudu}.

\paragraph{Současné trendy}
V současných trendech však dominuje konflikt na Ukrajině.
Je známé, že Rusko provádělo přípravu roky před samotným vypuknutím války na Ukrajině.
Zdroje poukazují na incidenty již v roce 2014\cite{russian_cyber_usage}.
Asi nejničivějším je však NotPetya z roku 2017\cite{NotPetya}
Faktem také je, že Rusko se nikdy neostýchalo využívat cyber útoky ve svůj prospěch.
Naopak, jejich používání je často těžkopádné a dalo by se říci i nezodpovědné.
Příkladem je útok zaměřený na olympijské hry ve městě Tokyo\cite{Tokyo_olympic_attack}.
Není tak překvapením zaznamenání hned několika cyber útoků před samotným vypuknutím války\cite{before_ukranian_war_attacks}.
Zajímavé jsou však hlavně markantní změny, které nastaly po samotném vypuknutí konfliktu a jejich vliv na bezpečnost do budoucna.
Možná nejlepším průvodcem těmito změnami je každoroční report vydávaný \firstac{ENISA}\cite{Enisa_thread_landscape}.
Z tohoto reportu jsou snadno rozpoznatelné hlavní trendy, stačí zohlednit celkové rozdělení reportu.
Hlavními zájmy jsou \textit{Ransomware ,Malware, Sociální Inženýrství, Hrozby pro data,	Hrozby pro dostupnost, Dezinformace a Supply-chain útoky}.
Samotné hrozby jsou blíže vysvětleny v následujících kapitolách.
Důležité je poukázat na trendů útoku na supply-chain a dezinformace.
Nelze říci, že tyto útoky byly v předchozích letech neexistující, v porovnání se současnou situací to tak ale téměř vypadá.
Jejich závažnost a také množství rapidně vzrostlo.
Není tak divu, že se objevují na listu hlavních hrozeb.
~\cntcapfigure{ENISA_prime_threats}{\linewidth}{Primární hrozby identifikované agenturou ENISA}{\cite{Enisa_thread_landscape}}\label{fig:PrimeThreats}

\paragraph{Pegasus spyware}
Speciální zmínku si zaslouží Pegasus spyware a celkově průmysl špehování na zakázku.
Za zrodem Pegasus spyware stojí Izraelské společnost NSO Group.
Pegasus je využíván pro špehování, překvapivě jsou však cílovým zákazníkem státní instituce a totalitní režimy.
Průmysl špehování na zakázku funguje celosvětově a zaměřuje se stále více na běžné uživatele internetu.
Cílem je nejen zisk informací ale také manipulace cílů za účelem kompromitování jejich zařízení a účtů.
Většina zprostředkovatelů těchto služeb tvrdí, že služby jsou určené pro sledování zločinců a teroristů.
Tomuto tvrzení však jednoznačně protiřečí  výzkum provedený společností \ac{ENISA}.
Více než 30,000 aktivistů za lidská práva, novinářů a právníků po celém světě bylo cílem špehování.
V Evropě byl hlášen vysoký počet případů právě se špionážním softwarem Pegasus od NSO.
Například byli monitorováni politici Španělska a vůdci katalánské nezávislosti.
Navíc sama skupina NSO uvedla, že pět států \firstac{EU} používá jejich špionážní software Pegasus.\cite{Enisa_thread_landscape}

\paragraph{Další entity v průmyslu špehování na zakázku}
~\cntcapfigure{Facebook_banned_entities}{\linewidth}{Entity identifikované společností Meta}{\cite{pegasus_facebook}}\label{fig:FacebookThreads}
Facebook je jeden ze zdrojů informací, které se používají pro špehování.
Meta se snaží tuto činnost zredukovat a tak provedla kroky proti 7 různým entitám operujícím špehování na zakázku.
Tyto entity sídlily působícím z Číny, Izraele, Indie a Severní Makedonie.
Entity zaměřovaly svoje působení na lidi ve více než 100 zemích po celém světě.
Obrázek ~\ref{fig:FacebookThreads} ukazuje ovlivněné entity.
Také je vidět, na kterých částech řetězce sledovací operace se jednotlivé entity podílely.
Řetězec se skládá z \textit{průzkumu}, \textit{navázání kontaktu} a \textit{využití}.
\textit{Průzkum} je první fáze řetězce sledování.
Obvykle je nejhůře zaznamenatelná.
Cíl sledování je potichu profilován s využitím software k automatickému sběru dat z celého internetu.
Informace o cíli pochází ze všech dostupných online zdrojů, a to jak volně dostupných tak i placených nebo nelegálních.
Jsou to data z veřejných webů, jako jsou blogy nebo sociální média.
Dále data z platforem pro správu znalostí, jako je například Wikipedia.
V neposlední řadě jsou data sbírána z "dark web" stránek, kde lze naleznout výpisy z napadených databází.
\textit{Navázání kontaktu} je druhá fáze řetězce.
Typicky je právě zde největší možnost odhalit sledování a s ním spojený útok.
Obezřetnost je tedy opět nejdůležitější krok k předcházení hrozeb.
tato část řetězce je zaměřena na navázání kontaktu s cíli nebo lidmi v jejich blízkosti.
Primárním cílem je vybudovat s cílem důvěru, získat nebo ověřit informace z předchozího kroku a v neposlední řadě připravit cíl na třetí fázi řetězce.
Za tímto účelem jsou zde využívány taktiky sociálního inženýrství.
Používají se fiktivní nebo ukradené identity aby oslovili lidi prostřednictvím e-mailů, telefonních hovorů, textových zpráv nebo zpráv na sociálních médiích.
\textit{Využitím} se myslí manipulace cíle.
Tato poslední část řetězce sledování obnáší vytváření phishingových domén, které jsou navržené tak, aby cíle oklamaly.
domený se snaží získat od cíle informace noho ho přesvědčit k provedení akce, která není v jeho nejlepším zájmu.
Jako cenné informace jsou brány přihlašovací údaje, například k  e-mailovému účtu, sociálním médiím, finančním službám nebo do  firemní sítě.
za nežádoucí akce se považuje instalace malware, specificky spyware pro další zisk informací nebo ransomware pro peněžitý zisk.\cite{pegasus_facebook}

\paragraph{Nejvíce napadené odvětví}
~\cntcapfigure{ENISA_targeted_sectors}{\linewidth}{Nejvíce napadené sektory (v období červenec 2021 - červen 2022)}{\cite{Enisa_thread_landscape}}\label{fig:TargetedSectors}
Z již zmíněného reportu\cite{Enisa_thread_landscape} je možné vyčíst, ve kterých odvětvích je zaznamenáno nejvíce útoků.
Útoky jsou jen zřídka omezené na jeden sektor, ve většině případů i útok zaměřený na specifický sektor přetéká díky provázanosti do ostatních.
Tento pohled na sektory je i přesto důležitý právě kvůli připravenosti jednotlivých sektorů na útoky.
Můžeme očekávat, že sektor zabývající se zprostředkováním digitálních služeb bude mít jinak připravenou ochranu než například vládní sektor.
A to nejen díky primárnímu zaměření sektoru ale také dopadu úspěšných útoků v minulosti a předvídaném dopadu budoucích útoků.
Na první pohled je z ~\ref{fig:TargetedSectors} vidět, že téměř čtvrtina útoků byla zaznamenána ve vládním sektoru.
Kolem 13 procent pak spadá jak sektoru zprostředkování digitálních služeb tak veřejnosti.
Lze očekávat, že množství útoků v těchto sektorech je ovlivněné dopadem, snadností propagace útoku a v případě veřejnosti taky celkovému množství možných cílů.
Zajímavá je také poznámka v reportu týkající se sektorů bankovnictví a zdravotnictví.
V době sledování byly tyto sektory napadané konzistentně v rámci celé sledované periody.
Zajímavý, hlavně v kontextu wálky na Ukrajině, je také armádní sektor, kde je zaznamenáno kolem 3 procent útoků.
Tato nízká procentní reprezentace by mohla být způsobena metodologií sestavování tohoto reportu.
Útoky byly zaznamenávány pomocí \firstac{OSINT}, neboli volně dostupných informací\cite{moje_bakalarka}.
Volně dostupné informace mají mnohá omezení.
Je jich velké množství a proto je nutné užitečná data dolovat.
Jen málokdy je na první pohled zřejmé, která informace je užitečná.
Hlavním problémem v tomto případě je však jejich vznik.
Volné informace jsou volné protože je tvůrce volně zpřístupnil široké veřejnosti.
I když společnosti začínají být zodpovědnější v nahlašování útoků, stále lze předpokládat že některé byly přehlédnuty nebo nenahlášeny.
Ve vojenském sektoru bude limitace informací prioritou.
Proto bude v kombinaci s propagandou obou stran konfliktu a limitací nežádoucích informací velmi obtížné sestavovat jakoukoli přesnou reprezentaci útoků.

\section{Nejčastější cyber-útoky}
Mezi nejčastější útoky dlouhodobě patří \textit{denial-of-service}, \textit{man-in-the-middle}, \textit{malware} a samozřejmě \textit{sociální inženýrství}\cite{LI20218176}.
Také zaznamenané útoky v roce 2022 jsou v souladu s tímto trendem.
Jsou to však velmi obecné kategorie, kde se specifika velmi liší.
Například pod malware se řadí již zmíněný spyware ale také ransomware, adware, rootkits, etc.
Pouze použití ransomware v roce 2022 vzrostlo o 41\%, primárně v první polovině roku.
Také čas na odhalení a nápravů úspěšného útoku trval o 49 dní déle, než bylo běžné.
V kategorii sociální inženýrství převládá phishing.
Ten je stále nejčastějším útokem.
Jeho oblíbenost spočívá v jeho efektivitě a relativně nízkých nákladech na provedení.
Počet zaznamenaných útoků téměř každoročně stoupá\cite{moje_bakalarka}.
Není tak divu, že jen za první polovinu roku 2022 byl zaznamenán nárůst v četnosti útoků o 48\%.
V této době bylo nahlášeno přes 11 tisíc útoků které stály cíle 12.3 milionů dolarů.
Je nutné poznamenat, že statistiky se týkají pouze \textit{zaznamenaných} útoků.
Mnoho společností tyto útoky veřejně neoznamuje.
Může tomu být z pokusu zachování dobré pověsti.
Horší možnost je že útok nezaznamenali.\cite{securityintelligence_most_common_2022_attacks}

\paragraph{Nové cíle}
Stále více útoků se zaměřuje na \textit{Supply-chain} a \ac{IoT}.
To nemůže být překvapením.
Již v předchozích letech bylo možné zaznamenat neklid mezi ruskem a ukrajinou.
Plné vypuknutí války samozřejmě zahrnuje útoky na infrastrukturu, a to i v online prostředí.
\textit{Supply-chain} je výborným cílem ze dvou důvodů.
Úspěšný útok v jedné části řetězce dovoluje šířit se po celku a dopad úspěšného útoku je proto významný.
Výzkumná zpráva od Accenture\cite{accenture_supply-chain-disruption} odhadla škody způsobené v Eurozóně těmito útoky na 112.7 bilionů euro.
Odhad na období 2022/2023 spojený právě s várkou je 242 bilionů euro a v nejhorším případě až 920 bilionů.
Také útoky na \ac{IoT} zařízení stoupají.
Není tomu však díky zvýšené oblíbenosti u útočníků nebo lepších metodách útoku.
Zvýšení plyne téměř čistě ze zvyšujícího množství zařízení.
\ac{IoT} zařízení jsou používána všude, najdeme je v domácnostech, firmách, továrnách nebo kasinech.
Tato zařízení ulehčují lidský život, bohužel jsou také novým rizikem.
Obzvláště protože jejich zabezpečení není priorita při vývoji těchto zařízení.
Proto jsou cílem útočníku at už jako přídavek do botnetu nebo jako slabý článek obrany samotné.
Zní opravdu špatně když se povede útočníkům ukrást důvěrné informace z kasína pomocí chytrého thermometeru akvária\cite{casino_thermometer}.
O vážnosti situace svědčí i veřejná prohlášení o slabinách, která vydávají přímo státy.
Například pro monitorovací zařízení vozidel MiCODUS MV720 toto prohlášení vydala americká vláda\cite{america_gov_advisories_micodus}.
A není divu, toto zařízení nejen monitoruje vozidlo ale je i schopné odpojit přívod paliva.
Jelikož je toto zařízení nachází hlavně v armádní a policejní technice, je to opravdu velký problém.\cite{securityintelligence_most_common_2022_attacks}

\paragraph{Zdroje nejčastějších útoků}
Právě z těchto důvodů je nutné podívat se nejčastější kategorie útoků.
V předchozím textu byly mnohé zmíněny, následuje jejich rozbor do hloubky.
Kategorie byly zvoleny z kombinaci několika zdrojů, jako je \ac{ENISA}\cite{Enisa_thread_landscape}, Cisco\cite{cisco_most_common_attack} nebo \firstac{NIST}\cite{nist_risks}.


%todo add definitions from att&ck
%todo add defending? prevention, good security higiene(segmentation, backups atd)
%todo add enisa -
hole\cite{Enisa_thread_landscape}

\subsection{Malware}
Tento termín se používá pro škodlivý software.
Zahrnuje známé hrozby jako viry, červy, trojské koně, spyware a ransomware.
Do systému se dostávají pomocí jeho slabin, často díky nepozornosti uživatele.\cite{cisco_most_common_attack}
Malware je důvod, proč si uživatele musí dávat pozor na jaké linky klikají a jaké programy instalují.
Je také nutné dávat si pozor odkud je program stažen.
Stále častěji malware předstírá, že je legitimní software.
Mezi nejznámější napodobované programy patří Skype, Adobe Acrobat, VLC nebo 7zip.\cite{virustotal_report_findings}
Škodlivému software se bohužel nevyhnul ani státní sektor.
Není však jen jeho cílem ale je také používán k šíření.
Zkoumání z roku 2022 ukazuje, že je možné najít desítky instancí hostování škodlivého software přímo v doménách spojených se státními institucemi.\cite{virustotal_report_findings_goverment}
Tento problém není specifický pro jeden stát ale je to celosvětový problém, jak je možné vidět na ~\ref{fig:MaliciousGoverments}.
Státní stránky nejsou blokované, naopak k nim existuje implicitní důvěra.
Proto je tento problém opravdu závažný.
~\cntcapfigure{malicious_govs_domains}{\linewidth}{Zneužití státní infrastruktury}{\cite{virustotal_report_findings_goverment}}\label{fig:MaliciousGoverments}

\paragraph{Obrana a detekce}
Obrana proti malware je složitá.
Nelze totiž ukázat na jednu specifickou chybu, kterou je nutné opravit aby tento problém nenastal znovu.
Existují však dobré uživatelské praktiky, které riziko infikování stroje malware snižují.
Absolutním základem je aktualizování softwaru používaného zařízení.
Mnoho velkých firem aktivně sleduje bezpečnostní incidenty spojené s jejich zařízeními a snaží se své uživatele chránit pomocí bezpečnostních aktualizací.
Pro detekci malware je vhodný antivirus, který je schopný z online databází nalézt nejznámější rysy malware kampaní a ty pak detekovat v nainstalovaných souborech.
Pokud nastane podezření na malware, je proto dobré pomocí antiviry zkontrolovat soubory.
Podezření by mělo nastat při výrazném zpomalení zařízení, nemožnosti vypínat programy, vyskakování reklam nebo cizích pokusech o přihlášení k osobním účtům.
Aby se předešlo samotné instalaci malware, je nutné dbát na bezpečné chování uživatele co se týče procházení internetu, instalace programů a klikání na soubory nebo odkazy v emailech.
Jak již bylo zmíněno, malware může být zabalen společně s legitimním software, je proto velmi důležité instalovat programy pouze z ověřených, ideálně oficiálních zdrojů.
Pozor na stránky, které napodobují oficiální stránky daného produktu.
Na nich lze malware očekávat s jistotou.
Nutné je také dávat pozor na podezřelé odkazy nebo soubory v emaily.
I přes tyto vektory se do zařízení může dostat malware.
V poslední řadě je nutné zmínit USB zařízení.
Pokud je to možné, uživatel by neměl nikdy zapojovat ke svému zařízení neznámé usb zařízení.
Ty jsou totiž nejen hrozba v podobě přenašeče malware ale dokonce i hrozba samotnému zařízení.
Příkladem takto nebezpečného zařízení je USBKill\cite{usb_kill}.
Toto zařízení na testování hardwaru se při napojení do usb portu nabije a pak vypustí svůj náboj zpět do portu.
Novější verze jsou dokonce schopné nabití dopředu.
To znamená, že jsou schopna útočit na vypnutá zařízení a obchází všechny známé bezpečnostní protokoly.

\subsubsection{Ransomware}
Za zvláštní zmínku stojí ransomware, jelikož je to velmi nebezpečný a útočníky oblíbený malware.
Jedná se o typ útoku, kdy útočníci přebírají kontrolu nad cílovými daty a požadují výkupné výměnou za obnovení dostupnosti těchto dat nebo zachování tajnosti dat.
Útočníci tak přímo ohrožují dostupnost, důvěrnost i integritu dat.
Jednoduché ransomware útoky data pouze odcizí a vyhrožují vypuštěním, popřípadě uzamknou systém, aby nebylo možné se k datům dostat.
Nejčastější jsou však sofistikované ransomware útoky, kde jsou data zašifrována a klič nebo dokonce program pro jejich odšifrování je oběti poskytnut až po zaplacení výkupného.
Ve výjimečných případech se stává, že klíč k odšifrování dat není poskytnut nikdy.
Důvod proč jsou tyto situace výjimečné je však zřejmý.
Není to dobré pro business ransomware skupin.
Nejen že neposkytování klíčů by vedlo k neplacení výkupného ale také to snižuje okruh cílů.
Zašifrovaná data totiž nemá cenu znovu napadat a žádat o další výkupné.
Tyto myšlenkové pochody útočníků nemůžou být překvapením.
Z interních dokumentů a logů skupiny Conti je zřejmé, že alespoň některé ransomware skupiny operují jako běžná společnost\cite{Conti_logs}.
Ve skupině pracovali manažeři, útočníci pracovali převážně během pracovního týdne a dokonce měli nárok na placenou dovolenou a jiné benefity.
K rozšíření tohoto typy útoku přispívá také koncept \ac{RaaS}.
Ten dovoluje i nezkušeným útočníkům operovat s velice sofistikovaným software.
Naopak zkušení útočníci nemusí ztrácet čas hledáním cílů ale soustředí se pouze na vyvíjení ransomware nástrojů a technickou podporu.
Z každého úspěšně provedeného útoku jim totiž náleží část peněz.\cite{Enisa_thread_landscape}


\paragraph{Monitorování}
~\cntcapfigure{ENISA_Ransomware}{\linewidth}{Počet zaznamenaných Ransomware incidentů a velikost ukradených dat ve sledovaném období květen 2021 až červen 2022}{\cite{Enisa_thread_landscape}}\label{fig:ENISA_Ransomware}
Jak je vidět z ~\ref{fig:ENISA_Ransomware}, Ransomware útoky byly v době sledování, tedy období od Května 2021 do červen 2022, časté a dopad na data byl značný.
Sledování těchto útoků je však velmi obtížné.
Pokud dochází pouze k vydírání, firmy mohou výkupné zaplatit a pokusit se zamlčet existenci útoku.
Se stoupajícími částkami výkupného je tato možnost stále méně validní.
Spíše se firmy obrací na pojištění proti ransomware.
To však ironicky zvyšuje jejich atraktivnost jako cíl, protože je jistota zaplacení výkupného
Dalším problémem monitorování je abstrakce ze strany útočníka.
Pro zmatení vyšetřovatelů nejsou pro ransomware skupiny neobvyklé změny názvu spolu s falešným odchodem do důchodu.
Tato stejná skupina pak se změněným jménem operuje nadále, můžou však pozměnit styl provádění útoků nebo nástroje.
Díky rivalitě mezi skupinami často dochází ke kradení zdrojových kódy, což však znesnadňuje identifikování této skupiny se skupinou původní.
Častým vodítkem bývají adresy krypto peněženek.
Výplaty výkupného v jsou vždy v kryptoměně, například Bitcoinu.
To samozřejmě také ztěžuje vystopování a trestní stíhání pachatelů.
Bez nadsázky lze říci, že monitorování ransomware skupin je složitější než se ransomware bránit.\cite{Enisa_thread_landscape}

\paragraph{Změny v metodice}
Ze zaznamenaných dat jasně vyplývá pokles útoky s počátečním vektorem prostřednictvím \firstac{RDP}, neboli protokolem vzdálené plochy.
Stále je to však druhý nejpoužívanější vektor útoku pro ransomware.
Útočníci stále prolamují slabá \ac{RDP} přihlašovací jména a hesla, zejména pokud není povolena \firstac{MFA}, v překladu vícefaktorová autentizace.
\ac{MFA} je velmi dobrou obranou proti neznalým útočníkům a jednodušším útoků, bohužel není nepřekonatelná.
Nejpoužívanějším vektorem útoku je samozřejmě phishing.
Za povšimnutí stojí nízké nakladu na využití těchto metod útoku a jejich značný výnos, což je jeden z hlavních důvodů jejich oblíbenosti mezi útočníky
Zajímavým trendem jsou změny ve způsobu vydírání společností.
Klasický útok ransomware nejdříve sbírá informace, poté k nim zamezí přístup a posléze vydírá společnost.
Teprve pokud společnost odmítne zaplatit výkupné je útok veřejně oznámen.
Získaná data jsou také prodávána nebo volně dostupná na internetových fórech dark webu.
Některé ransomware skupiny však začali používat jiný postup.
Oběti je nabídnuto, aby si data \textit{koupila} dříve, než je zakoupí její konkurenti.
Útočníci také již po provedení útoku změnili vzhled webových stránek oběti aby o útoku oznámili.
Takto prováděný útok vede k veřejnému hanobení oběti útoku a má přímý vliv na vztahy s veřejností a reputaci společnosti.
Dalším znepokojivým trendem je vytváření dedikovaných stránek s daty o napadených přímo na veřejném internetu.
Tento postup zvolil ALPHV, také známý jako BlackCat ransomware, který operuje jako \ac{RaaS}\cite{ALPHV_blackcat}.
Na tomto webu se nacházely informace o zákaznících i zaměstnancích a ti si tak mohli snadno zkontrolovat, zda byli útokem zasaženi.
Vytvoření takto veřejného repozitáře se stalo novým způsobem, jak donutit třetí strany, aby se identifikovaly.
Lze předpokládat, že postižení zákazníci sami kontaktují majitele webu aby zaplatili za stažení dat.

\subsubsection{Spyware}
%todo add


\subsection{Sociální inženýrství}

%todo from my bakalářka

Phishing útoky jsou prováděny ve velkém objemu a cílí na širokou veřejnost, zatímco jiné útoky sociálního inženýrství využívají vlastní kampaně přizpůsobené konkrétním zaměstnancům. Při použití sociálního inženýrství útočníci využívají přístup zaměstnanců uvnitř organizace k získání technického základu v síti, ze které provádějí další útoky.

\subsubsection{Phishing}
Podvodné komunikaci, která předstírá že pochází z legitimního zdroje, se říká Phishing\cite{cisco_most_common_attack}.
Nejčastěji se jedná o emailovou komunikaci, existují však i formy přes sms nebo telefonní hovor.
Nejnebezpečnější jsou specializované útoky, zvané whaling a spear-phishing.
Většina útoků se snaží pokrýt velké spektrum populace, většinou omezené pouze záminkou, která je k útoku využita.
Specializované útoky jsou však tvořeny pro jednu specifickou osoby.
Útočník stráví týdny nebo i měsíce zjišťováním informací o cílené osobě a poté vytváří jeden specifický útok, který je v drtivé většině případů úspěšný.
Tato hrozba mívá 2 primární cíle.
Prvním je ověření nebo zisk informací, které lze použít k dalšímu útoku.
Druhým je kompromitace zařízení, například pomocí malware.
Nejčastěji je to pomocí odkazů, které se v podvodném emailu nachází, popřípadě pomocí infikovaných souborů.
Znepokojující je také každoroční nárostu počtu těchto útoků.\cite{moje_bakalarka}




\subsection{Threat against data}

\subsection{Sql injection}
Tento známý útok využívá špatného nastavení databáze pro provádění za normálních okolností nepovolených operací.
Ze špatně zabezpečeného vyhledávacího boxu je útočných schopný zadáním specificky vypadajícího dotazu získat například výpis celé databáze.
%todo more? zdroje?

\subsection{Threat against availability}

\subsubsection{Denial of service}
Útok zaměřený na kolaps infrastruktury.
\textit{Denial-of-service}, také znám jako \firstac{DoS} útok se snaží vyčerpat zdroje služby množství zasílaných požadavků.
Další možností je ovládnutí nebo alespoň vyřazení z provozu dané služby za pomocí známe zranitelnosti.\cite{cisco_most_common_attack}


Známější a více používaná je však podkategorie tohoto útoky zvaná \firstac{DDoS}.
Hlavní změnou je využívání mnoha zařízení, takzvaného \textit{Botnetu}, která na příkaz zahltí službu množstvím požadavků.
Problematika tohoto útoku spočívá v 2 hlavních bodech.
Při útoku \ac{DoS} je možné příchozí požadavky z jednoho zařízení snadno zablokovat pomocí \textit{blacklistu} daného zařízení.
To však není proveditelné pokud tyto požadavky přichází z tisíců nebo i statisíců zařízení.
S tím souvisí i druhý problém, rozsah celého útoku.
Jedno zařízení jen těžce může konkurovat statisícím v množství požadavků, které je schopné odeslat.
A to i za předpokladu, že je mnohem výkonnější.
%todo more? zdroje?

\subsubsection{Internet threats}


\subsection{Ostatní}
Tato kategorie se skládá z útoků, hrozeb a technik, které nelze zařadit samostatně do jedné kategorie ale jsou často využívány v kombinaci s jinými kategoriemi útoků.
je nutné je v tomto výčtu zmínit kvůli jejich nebezpečnosti, rozšířenosti a nízkému veřejnému povědomí.

\subsubsection{Man-in-the-middle}
Útoky \textit{Man-in-the-middle} jsou také zvané odposlouchávací útoky.
Útočník se při tomto útoku dostává doprostřed komunikace 2 subjektů.
Hlavním cílem je filtrování, kradení a nahrazování dat proudících v komunikaci.
Mezi nejčastější implementace tohoto útoky patří napodobování veřejné \textit{Wi-Fi}.
Útočník se představuje jako vstupní bod na tuto veřejnou wifi kap posléze odesílá veškerou komunikaci.
Lidé si tak nemusí uvědomit, že útočník existuje a nedávají si tak pozor na svoji činnost na takto odposlouchávané síti.
Další oblíbenou implementací útoku je použití malware.
Ten po nainstalování přesměrovává komunikaci nebo pouze odposlouchává činnost na daném stroji.\cite{cisco_most_common_attack}
%todo more?

\subsubsection{Zero day exploit}
\textit{Zero day exploit} je slabina na kterou ještě neexistuje patch.
Tato situace může nastat pokud byla slabina nahlášena a ještě ji výrobce systému nestihl opravit.
Mnohem častější je případ, kdy slabinu našel sám útočník.
Ten ji pouze bude chtít využít, rozhodně ne nahlásit.
Tyto slabiny se také snaží získat státní složky zaměřené na špehování.
Nemůže tak být divu, že s těmito slabinami existuje černý trh\cite{world_end_2021} kde se prodávají až za miliony dolarů.
%todo more? zdroje?

\subsubsection{DNS tunneling}
\textit{DNS tunneling} je možné použít pro legitimní účely, tato práce však zkoumá jeho zneužití.
Není to samotný útok ale podpora útoku.
Je založený na \firstac{DNS}.
Tento protokol je využíván prakticky neustále pro běžné účely a je velmi flexibilní.
To však dovoluje jeho zneužití.
\ac{DNS} je tak často používán, že malé navýšení není poznat.
Mnoho společností ho navíc nekontroluje dostatečně, aby mohlo tento typ útoku možné odhalit.
Útočník v \ac{DNS} requestech a responsech zakódovává vlastní komunikaci za účelem jejího skrytí.
Odchozí komunikace může být použita pro exfiltraci kritických dat nebo přijímání odpovědí od nainstalovaného malware.
příchozí naopak pro posílání příkazů již nainstalovanému malware nebo dokonce stahování dalšího malware.\cite{cisco_most_common_attack, bright_sec_dns_tunneling}
%todo more? zdroje?

\subsubsection{Supply-chain útoky}
%todo add

\section{Kategorie útočníků}
Pro kompletnost je při zkoumání bezpečnostního nutné uvažovat i typu útočníka, který útok uskutečňuje.
Různé typy se chovají během útoku různě, mají jiné cíle, možnosti a motivace.
Pokud je cílem komplexní ochrana, je nutné zvážit, se kterými typy útočníku bude nejčastější konflikt.
Pro zjednodušení je možné útočníky zařadit do několika skupin.
Stejně jako u útoků mezi skupinami existují návaznosti a překrývání, je nutné skupiny brát pouze jako hrubé zařazení a ne definitivní výčet.
Skupinami jsou:
\begin{itemize}
	\item Příležitostný útočník,
	\item Script kiddie,
	\item Hacktivista,
	\item Hacker na objednávku,
	\item Cyber kriminálník,
	\item Státem podporovaný útočník,
	\item Insider (vnitřní osoba).
\end{itemize}\cite{Enisa_thread_landscape,data_flair_attackers}

\subsection{Příležitostný útočník}
Tento typ útočníka by bylo možné nazývat také hobby útočník.
Mají omezené zdroje a bývají velmi specifičtí v rámci postupů a zranitelností, které využívají k dosažení svého záměru.
Tito útočníci nejsou plně zaměřeni na hackování, je to pro ně spíše zábava a zpestření života.
Z toho také vyplývají hlavní motivace těchto útočníků.
Cílem je získat proslulost.
Existují fora dedikované těmto útočníkům, kde se navzájem chlubí svými útoky ostatním.
Můžeme tak očekávat, že útoky budou směřovány k jasně viditelným avšak málo destruktivním typům.

\subsection{Script kiddie}
Vě většině případů nejméně schopní útočníci.
Jsou to amatéři, kteří používají volně dostupné nebo zakoupené nástroje bez hlubších znalostí systémů a hackování.
Většinou bývají neškodní, teprve se učí a mají omezené znalosti a možnosti útoku.
Proto jsou jejich útoky málokdy úspěšné a i v případě úspěchu často nevědí, jak pokračovat.
Lze tak předpokládat nízkou úroveň poškození u útoků, snadno odhalitelné útoky a nízkou sofistikovanost.
Jediným výrazným rizikem jsou zakoupené nástroje pro specifické účely.
Ve výjimečných případech kdy bývá motivací pomsta zakupují script kiddies na černém trhu nástroje pro konkrétní devastující útok.
Ve většině případů se však jedná o nováčky kteří hledají výzvy a adrenalin.
Pokud má skript kiddie pro hackování talent, může se díky zisku zkušeností a znalostí stát v průběhu času profesionálním hackerem.

\subsection{Hacktivista}
Útočníci motivování politickými nebo ideologickými názory jsou nazýváni hackticvisti.
Provádějí škodlivé útoky pro zviditelnění a prosazení své agendy.
Možná nejznámějším příslušníkem je skupina Anonymous.
S velikostí skupiny je očekávaná jistá fragmentace, a tak se anonymous hlásí k množství incidentů s různými agendami.
Za speciální zmínku stojí jejich kybernetické útoky proti Rusku v souvislosti s válkou na ukrajině\cite{anonymous_russian_attacks}.
Jak je jasné, motivací hacktivistů není zisk, ale snaha bojovat za spravedlnost.
Jejich útoky tak bývají ve vážnosti v rozmezí od neškodných, jako je změna webové stránky, po drastické kde vede k poškození infrastruktury a zničení nebo vyzrazení citlivých dat.
Také není neobvyklá spolupráce s Insidery, kteří sdílí jejich přesvědčení nebo je upozornili na objekt jejich zájmu.
Lze tedy říci, že hacktivisté mohou dle svého přesvědčení být kvalifikování jako síla dobra i zla.


\subsection{Hacker na objednávku}
Skupina schopných hackerů kteří své vědomosti a znalosti propůjčují jiným osobám či organizacím jsou hackeři na objednávku.
Nemůže být překvapením, že velká většina hackerů více či méně spadá do této skupiny.
Skupina se totiž navíc dle motivace,metodologie a záměru dělí na podskupiny, zvané \textit{White hat, Gray hat} a \textit{Black hat} hackeři.

\paragraph{}
\textit{White hat hackeři} využívají své schopnosti ve spolupráci s organizacemi k určení jejich slabých míst v zabezpečení.
Tito hackeři se řídí silným etickým kompasem, který je jasně zřetelný z jejich činnosti a zaměření pro zlepšení zabezpečení.
Pracují s svolením organizace a v jasně vymezených hranicích, což však může vázat jejich schopnost odhalit veškeré slabiny.

\paragraph{}
\textit{Black hat hackeři} jsou neetický opak white hat hackerů.
Jejich primární metodikou je prolamování obrany systémů a sítí bez vědomí organizace.
Cílem těchto hackerů je vždy svůj vlastní prospěch.
Existuje mnoho možností, jak tohoto prospěchu dosáhnou.
Někteří napadají společnosti a exfiltrují data uživatelů, která pak prodávají na černém trhu.
Jiný využívají tato data k bankovním podvodům, kampaním sociálního inženýrství nebo kradení účtů.
Další možností je pouze vyhledávání slabin a jejich následný prodej na trhu slabin.
Státní organizace jsou za takzvané \textit{zero day slabiny} ochotné zaplatit i statisíce dolarů\cite{world_end_2021}.
Tam kde by slabinu etický hacker ohlásil aby mohla být napravena ji neetický hacker prodá za značnou sumu peněz.

\paragraph{}
\textit{Gray hat hackeři} jsou na pomezí předchozích kategorií.
Jestli boudou klasifikování blíže k white hat hackerovi nebo black hat hackerovi závisí na situaci.
Mnohdy využívají metodologii black hat hackerů k odhalení slabin, ale místo prodeje nebo zneužití slabiny ji dané organizaci nahlásí.
Nesnaží se tak způsobovat škodu, ale odhalit slabiny zabezpečení, které by mohly být nemožné odhalit pro white hat hackera.
To však bohužel stále znamená napadení organizace či osoby a ty nemusí reagovat kladně na pozdější oznámení útoku.
Také občasné požadavky platby za napravení slabiny nebo její nevyzrazení jim nedělá nejlepší jméno.
Asi nejlepším příkladem tohoto fenoménu je kauza okolo Twitteru Donalda Trumpa, jeho hesel a hackerů nazývaných \textit{The Guild of the Grumpy Old Hackers}\cite{darknet_87_grumpy}.
Této skupině hackerů se povedlo několikrát získat přístup k Twitter účtu Donalda Trumpa, a to v letech 2016 a 2020.
Pro kontext, Donalt Trump byl v roce 2016 kandidát na prezidenta Ameriky a v roce 2020 dokonce prezident.
Po obou těchto úspěšných útocích to sami hackeři nahlásili, avšak i tak se nevyhnuli soudnímu řízení.
Naštěstí byl jejich postup zcela legální a byly tak zproštěni viny.

\subsection{Cyber kriminálník}
Tato skupina útočníků jsou jednotlivci nebo dokonce organizované skupiny zaměřené na finanční zisk pomocí kriminální aktivity.
Mezi jejich činnosti patří ransomware, těžba kryptoměn, krádež kryptoměn nebo krádež přihlašovacích údajů, které využívají pro další činnosti.
Množství těchto útočníků stoupá také díky stále nižší vstupním nárokům díky \ac{MaaS}\cite{MaaS_ATaT} a \ac{RaaS}\cite{RaaS_microsoft}.
Zajímavým trendem posledních let je navíc neustálé přejmenovávání a `odcházení do důchodu' kriminálních skupin.
Cílem je vyhýbání se zákonnému stíhání a sankcím, jelikož je na tyto skupiny kladena stále větší pozornost.

\subsection{Státem podporovaný útočník}
Útočníci úzce spojeni s určitým státem.
Jejich motivace a cíle jsou vždy v zájmu jejich státu.
Útoky jsou tedy často zaměřené na posílení svého státu v politické nebo vojenské sféře.
Je jasné, že tito útočníci mají téměř neomezené zdroje, vyspělé technologie a prvotřídní znalosti.
Proto jsou od nich časté sofistikované a rozsáhlé útoky.
Lze tedy očekávat průmyslovou špionáž, úniky státních tajemství, cílené kybernetické útoky na infrastrukturu ale také rozsáhlé ransomware kampaně a podvody s kryptoměnou.
Není neobvyklé, že státy rekrutují úspěšné hackery ze všech možných zdrojů.
Existují spekulace o rekrutování i mezi kryminálníki.
Mnoho spekulací koluje například kolem skupiny REvil, kterou Rusko zatklo na začátku roku 2022\cite{REvil_story}.
Bohužel již o pár měsíců později se znovu začaly objevovat známky jejich aktivity.
Jak je však u státem sponzorovaných akcí trendem, daný stát jakoukoli vazbu popírá a tento fakt lze jen těžko vyvrátit.
Bylo identifikováno hned několik skupin, kde je návaznost na stát téměř jistá.
Příkladem je \textit{Lazarus Group} pro Severní koreu, \textit{Sandworm Team} pro Rusko nebo \textit{APT41} spojený s Čínou\cite{Mitre_groups}.

\subsection{Insider (vnitřní osoba)}
Insider je osoba pocházející z organizace, která způsobila chtěně nebo nechtěně bezpečnostní incident.
Představuje tak značnou hrozbu pro bezpečnost celé organizace.
Většina bezpečnosti je nastavena na venkovní hrozby.
Nebezpečí insidera je právě v naprosté neefektivitě takto postavené obrany vůči němu.
Systém nemůže být nikdy stejně chráněný z venku jako zevnitř
Insider takto nastavenou ochranu obejde a dostane se přímo ke zranitelným místům.
Velmi často jsou těmito insidery nespokojení zaměstnanci, bývalí zaměstnanci, dočasní pracovníci nebo dokonce zákazníci.
Je proto velmi důležité správně nastavit přístupy do systému a při odchodu zaměstnanců jim přístup včasně odebrat.
Insidery můžeme dle úmyslu dělit na tři hlavní kategorie, kterými jsou \textit{Úmyslný insider}, \textit{Nechtěný insider} a \textit{Nedbalý insider}.

\paragraph{Úmyslný insider}
Úmyslný insider je ten, kdo vědomě způsobí svojí organizaci újmu.
Mezi příklady se řadí krádeže dat, úmyslné poškození infrastruktury nebo vydání přístupu třetí neoprávněné osobě.
Cílem těchto insiderů bývá často pomsta nebo jen peněžní zisk.
Speciálním případem této kategorie je takzvaný \textit{whistleblowing}.
Tento termín označuje veřejné upozorňování na trestné, neetické nebo podezřelé jednání firem či státních institucí jejich zaměstnanci.
Tohoto zaměstnance lze tedy také považovat za úmyslného insidera, který však vyzrazuje informace v dobré víře a nejedná ve vlastní prospěch.
Velice známým zástupcem je Edward Snowden, který v roce 2013 veřejnosti odhalil existenci rozsáhlých programů na sběr informací fungujících pod záštitou \firstac{NSA}\cite{Edward_Snowden}.

\paragraph{Nechtěný insider}
Může se stát, že insiderem se stane běžný zaměstnanec.
Shodou okolností se dopustil nějaké chyby která poškodila firmu.
To však nebylo úmyslem a ani tato situace nenastala v rámci zanedbání povinností.
Může se například jednat o náhodné smazání nebo modifikace důležitých souborů, nechtěné vypojení zařízení ze zásuvky.

\paragraph{Nedbalý insider}
Některé hrozby vznikly kvůli nedbalosti zaměstnance.
Tne pak figuruje jako nedbalý insider.
Situace nastane, když zaměstnanec nedodržuje správné postupy, politiky nebo metodologie, které organizace stanovila.
Tyto postupy vzniklý právě pro ochranu důvěrnosti, integrity a dostupnost dat.
Jejich nedodržování je proto značný bezpečnostní problém.
Je však nutné zkoumat, zda toto nedodržování plyne z nerealisticky nastavených politik nebo pouze chyb zaměstnance.
Další možností jsou také nedostatečné znalosti a zaškolování zaměstnanců pro vykonávání možné činnosti.
Je jasné, že v mnoha případech může být obtížné rozeznat, zda se jedná o nedbalého nebo pouze o nechtěného insidera


\section{Testování}
Smyslem testování je simulovat útok na infrastrukturu systému.
Vě většině případů to znamená, že se specialista na bezpečnost snaží provádět stejné kroky jako potencionální útočník.
Většina případů však nejsou všechny.
Ideální postup vylučuje nejen nepředvídatelnost útočníka, ale ve výjimečných případech i jeho neschopnost nebo neznalost.
Stejně jako v každém odvětví existují i mezi útočníky sféry kompetence a znalostí.
Existuje mnoho útočníků, kteří se touto činností nejen živí ale je to také jejich vášeň.
Takovýto útočník bude disponovat znalostmi a technikami které nejsou běžné.
Při útoku bude postupovat co možná nejoptimálněji za účelem dosažení výsledků.
Je možné od něj očekávat pozorné zkoumání cíle a vytrvalost, díky které se si cíl zvolí.
Tiše hledá slabiny a proniká do střežených systémů bez povšimnutí.
Ne každý útočník je však takto nebezpečný.
Drtivá většina útoků je prováděna nadšenci nebo útočníky bez obsáhlých znalostí a technik útoku.
Není tak vyloučeno použití snadno detekovatelných `hlučných' technik nebo prostých chyb, kterými na sebe při útoku upozorní.
Nebezpečnost této skupiny v posledních letech vzrostla díky propagaci \firstac{MaaS}\cite{MaaS_ATaT} a \firstac{RaaS}\cite{RaaS_microsoft}.
Existence nástrojů, které dovolují provézt úspěšný útok za pomocí znalostí zkušenějšího útočníka znamená, že ani nezkušené útočníky není možné podceňovat.

\paragraph{}
Testování je jedna z nejlepších možností, jak zlepšit obranyschopnost organizace.
Nemilou samozřejmostí je, že největší změny v ochraně přichází po bezpečnostním incidentu.
Ztráta dat, poškození reputace organizace nebo i finanční postihy jsou výborné motivace nejen pro obránce ale i management organizace věnovat zabezpečení větší úsilí.
Pokud se útočníkovi podařilo prolomit ochranu organizace, jasně to ukazuje na nedostatky, které je nutné napravit.
Bezpečnostní incident tak se kromě mnoha nepříjemností ukazuje, jak dobře byla organizace připravena.
Dle reakce obránců lze zkoumat jejich celková připravenost.
Dobře připravení obránci mohou nebezpečný incident zaznamenat včas a pokud jeho následku nejde zabránit, mohou ho alespoň omezit.
Nepřipravení obránci nemusejí zaznamenat incident ani měsíce po jeho uskutečnění a několikanásobném zopakování.
Organizace a její obránci musí vědět, jak postupovat v případě incidentu.
Testování je simulace bezpečnostního incidentu ale bez většiny jeho negativ.
V závislosti na druhu testování je možné vyzkoušet připravenost obránců, jejich obezřetnost a postupy v případě detekce podezřelé činnosti.
Můžeme testovat robustnost infrastruktury, zda je kvalitně navržena a správně nastavena.
Můžeme testovat zařízení, zda fungují správně a nelze je zneužít.
Můžeme testovat technologie, zda jsou správně využívány popřípadě zda je jejich využívání nutno omezit nebo naopak rozšířit.
Důvodů, proč by mělo být testování součástí zlepšování obranyschopnosti každé organizace je tedy mnoho.
Testování přináší tolik benefitů protože to není prostá reakce na událost ale příprava.
Cílem organizace je zajistit, aby žádný bezpečnostní incident nebylo možné stejně uskutečnit dvakrát.
Čím více je poprvé uskutečněno pomocí testování, tím více jich bude zaznamenáno a zastaveno při skutečném útoku.


\subsection{Rozsah testovní}
Obránci se musí přizpůsobovat a testovat nejen nejběžnější a nejefektivnější útoky ale i krajní scénáře.
Toto přizpůsobování vedlo nejen k různým testovacím metodikám jako \textit{Red teaming} a \textit{Penetration testing} ale také k testování různých scénářů a hlavně rozsahů.
Mezi nejznámější patří \textit{celkový pokus o kompromitaci}, \textit{specifický pokus o kompromitaci} nebo \textit{předpokládaný kompromis}.
Je samozřejmé, že každá tato kategorie obsahuje další podkategorie.
Toto základní rozdělení je však vhodné pro rozpoznání hlavních výzev, složitostí ale i přínosů testování.
Nelze říci, zda je jeden scénář lepší nežli jiné, naopak se doplňují a každý slouží pro určitou situaci.\cite{red_team_oakley_2019}

\paragraph{Celkový pokus o kompromitaci}
Základní situací je \textit{celkový pokus o kompromitaci obrany}.
Lze definovat jako snaha \textit{red teamu} napadnout celou infrastrukturu organizace a pokusit se kompromitovat co možná největší možnou část.
Tato strategie je vhodná pro celkové testování infrastruktury.
Je to nejbližší emulace skutečného útoku.
Testování ve většině případů začíná ze stejného bodu, který je běžně dostupný veřejnosti a tím pádem i útočníkům.
Lze takto testovat správnou segmentaci infrastruktury a resistivity jednotlivých segmentů.
Z velkého a otevřeného rozsahu však vyplývá jen omezená možnost testování do hloubky a zaměření se na specifické části infrastruktury.
\textit{Red team} má pouze omezený čas a prostředky, proto nachází pouze nejčastější slabiny na celém rozsahu testování ale nemá možnost testovat jednotlivé části do hloubky a pro všechny možné hrozby.
Při nesprávném provedení nebo nepochopení cíle tohoto testování může vznikat falešný pocit bezpečnosti, který je pro zabezpečení nebezpečný.
cílem tohoto testu je získat celkové povědomí o ochraně a poté pokračovat s podrobnějším testováním.
%todo add picture?

\paragraph{Specifický pokus o kompromitaci}
Pro testování omezené části infrastruktury se využívá \textit{specifický pokus o kompromitaci}.
Cílem je důkladně otestovat omezenou část infrastruktury bez využití, zásahu nebo omezení částí infrastruktury nacházející se mimo testovaný rozsah.
\textit{Red team} prování podrobné testování, které při korektním provedení vede k nalezení všech slabin.
Výhodou tohoto testování je jeho specificita, je možné testovat určitou oblast infrastruktury nebo jen určité zařízení a to pouze za určitých podmínek.
Příkladem může být testování za použití specifického účtu a s ním spojených pravomocí přístupu.
Další výhoda je možnost testování v běžném provozu.
Krátkodobě vyřadit a testovat část infrastruktury je proveditelnější nežli celkové vyřazení nebo testování za běžného provozu.
V takovém případě je však nutné dávat pozor na správné vybrání rozsahu.
Při důkladném testování není neobvyklé vyřazení testovacího zařízení z běžného provozu nebo jeho zhroucení.
Je tak důležité aby nebyl testován systém, který je kriticky důležitý nebo zrovna používaný.

\paragraph{Předpokládaná kompromitace}
Z názvu je jasné, že toto testování zkoumá, jakou škodu může útočník napáchat při úspěšném útoku.
Počáteční bod tohoto testování je založen na předpokladu úspěšné infiltrace útočníkem.
Přístup k takto identifikovanému bodu je s příslušnými právy poskytnut \textit{red teamu}, který nemusí prolamovat obranu systému.
\textit{Red team} se tak může plně soustředit na zkoumání dopadu úspěšného útoku.
Důvod, proč toto testování není pouze součástí předchozích testování je úspora času a prostředků.
Vhodným příkladem jsou útoky sociálního inženýrství a phishing.
Jsou to často testované útoky díky jejich četnosti, jejich výsledek a prolomení systému však není okamžité.
Může trvat i týdny než uživatel otevře email a nainstaluje malware.\cite{red_team_oakley_2019}
Z důvodu efektivity je tedy lepší předpokládat úspěch útoku a začít testovat z bodu malware nakaženého uživatelského stroje.
Tento typ testování je výhodný pro ukázku dopadu úspěšného útoku.
Častým problémem bezpečnosti je její nedocenění dokud funguje správně.
Ve chvíli kdy nastane problém je však většinou pozdě.
Proto tento typ testování představuje nástroj prokázání důležitosti kvalitní obrany ukázkou důsledků jejího selhání.

\subsection{Metodologie testování}
ještě důležitějším rozhodnutím než jak široce testovat je způsob testování.
možná nejznámější jsou 3 hlavní testovací metodologie: Penetration testing, Red team assessment a Purple team/kontrolní cvičení.
Samozřejmě stejně jako i u rozsahu testování i u těchto metodologií existuje určitý přesah.
Z tohoto důvodu jsou tyto metodologie nepřesně definovány, někdy zaměňovány nebo dokonce považovány za jedinou metodologii s několik různými názvy.
Je proto nutné si je představit v kontextu tohoto textu.
K tomu slouží následující sekce.\cite{securityInteligence_pen_test_red_team_purple_team}

\paragraph{Penetration testing}

Penetration testing, často také zkracován na pen testing, je metodologie s cílem vyhledání co nejvíce možných zranitelností testovaného systému.
Takto důkladné hledání je omezeno pouze na zadaný `scope', neboli část testovaného systému.
Aby bylo dosaženo co největší efektivity, tento typ testování se svoji činnost nesnaží schovávat, naopak bývá velmi snadné ho detekovat, to však nevadí protože ochrana daného systému by o jeho provádění měla být předem informována.
Tento typ testování může být také z části automatizován, což vede k dalšímu zlepšení efektivnosti na úkor flexibilitě.
V neposlední řadě Penetration testing vyžaduje menší množství zdrojů, ať už v podobě času, lidí či kapitálu, než red team assessment.
\paragraph{}
Ve shrnutí, Penetration testing je hledání co nejvíce zranitelnosti v zadané oblasti bez nutnosti svoji činnost maskovat před ochranou systému.\cite{securityInteligence_pen_test_red_team_purple_team, lootsec_pen_test_vs_red_team, astra_pen_test_vs_red_teaming}


\paragraph{Red team assessment}

Red team assessment, také referovaný jako red teaming, je na rozdíl od pen-testingu zaměřen na dosažení zadaných cílů, například získání přístupu k citlivým datům, nebo testování všímavosti ochranných systémů.
To také znamená, že red teaming není omezen oblastí, kterou může v rámci dosažení svého cíle využít.
Obrana systému není ze zřejmých důvodů seznámena s nadcházejícím testováním, proto tato metodika vyžaduje maskování a opatrný přístup, aby pokud možno co nejvíce simulovala skutečný útok.
Výhodou je, že za předpokladu detekce testování bude obrana systému postupovat stejně, jako kdyby se jednalo o skutečnou hrozbu a získáme tak cenná data a zpětnou vazbu pro zlepšení budoucí obrany.
Je jasné, že pro správné provedení red teamingu je potřeba více zdrojů, at už se jedná o čas nebo lidskou práci, než v případě pen testingu.
Také automatizace zde bývá složitá, jedná se totiž o specifickou strategii pro každé testování kde je nutnost reagovat na všechny události vzniklé testováním.\cite{securityInteligence_pen_test_red_team_purple_team, lootsec_pen_test_vs_red_team, astra_pen_test_vs_red_teaming}

\paragraph{}
Ve shrnutí, red team assessment je zaměřen na testování obranyschopnosti cíle s omezeními vznikajícími pouze ze snahy simulovat reálný útok.


\paragraph{Purple team/Kontrolní cvičení}
Specifická situace je pak blízká spolupráce útočníků (Red team) a obránců (Blue Team).
Hlavní výhodou tohoto postupu je zpětná vazba mezi útočníky a obránci, kde se dá testovat detekce a následná odezva obránců při testování specifických technik obsažených v MITRE ATT\&CK frameworku.
Další výhodou je seznámení s postupy druhé strany pro tvorbu efektivnějším strategií, at už se jedná o útočníky nebo obránce.
\paragraph{}
Tato cvičení lze uskutečňovat pouze jednou pro vyzkoušení techniky nebo také opakovaně.
Částečná automatizace za použití open-source nástrojů, které mají v této práci vlastní kapitolu, je také jednou z možností.
Výhodou opakovaného testování s automatizací je možnost měřit a porovnávat výsledky za stejných stupních podmínek, což vede k jasně viditelnému zlepšení obranné strategie.\cite{securityInteligence_pen_test_red_team_purple_team,redscan_team_purple_team}

\section{Zavedené nástroje v problematice Testování}
Pro účely testování je dobré využívat stávající technologie.
Je jasné, že každé testování prostředí je jedinečné.
To však neznamená, že není možné těžit z poznatku ostatních obránců.
Vzniklo mnoho nástrojů, jejichž cílem je pomoci nejen s rozeznáváním útoků ale i se samotným testováním.
Asi nejznámějším frameworkem pro kolekci útoků a jejich rozeznávání je MITRE ATT\&CK\cite{mitre_attack_framework}.
Tento framework je open-source, volně dostupný a obsáhlý.
Existuje množství testovacích frameworků, které jsou s MITRE ATT\&CK kompatibilní\cite{csoonline_4_testing_frameworks}.
Příkladem \firstac{RTA}, Caldera nebo Atomic Red.
Tyto frameworky jsou na rozdíl od jiných zdarma a proto se hodí pro začátek testování.
Je dobré si nejdříve vymezit testovanou oblast.
Všechny tyto frameworky budou více probrány a porovnány v následujících kapitolách.

\subsubsection{MITRE ATT\&CK}
MITRE ATT\&CK je bezplatná otevřená znalostní databáze popisující chování útočníků a techniky, které tyto útočníci využívají.
Tato znalostní báze je založená na pozorování v rámci reálného světa, data pochází z veřejného zpravodajství a hlášených bezpečnostních incidentů.
Samozřejmě jsou začleněny i veřejné výzkumy technik útočníků.
Je aktivně udržována, nové informace jsou přidávány dvakrát do roka\cite{mitre_attack_framework_faq}.
Hlavním přínosem tohoto projektu však není pouze mapování útočníků, ale celková klasifikace útoků do jednotlivých \textit{taktik, technik a sub-technik}.
framework tak obráncům nabízí možnost rozložit složitý útok na jeho jednotlivé stavební kameny.
ATT\&CK je velmi podrobný a dokonce u jednotlivých technik útočníků nabízí možnosti detekce a obrany.
Obránci proto mohou snáze porozumět stylu práce útočníků, technikám které používají ale nejdůležitěji, jak tyto informace využít k vytvoření lepší obrany.
Všechny tyto informace jsou navíc přehledně dostupné v interaktivní formě pomocí ATT\&CK navigátoru\cite{mitre_attack_navigator}.
V něm je možné provádět vizualizaci a porovnávání útoků.
Je tedy snadné prezentovat rozbor útoku dalším osobám a díky exportování a importování je snadné studovat již popsané útoky a sledovat nejpoužívanější taktiky a techniky.
Není proto divu, že je MITRE ATT\&CK v době psaní tohoto textu považován jako jeden z nejvýznamnějších standardů pro popis technik používaných útočníky při kybernetických útocích.\cite{mitre_attack_framework}

%todo add picture of navigator

\subsubsection{Členění}
Díky popularitě ATT\&CK došlo k jeho velkému růstu.
Je jasné, že s přidáváním taktik a technik vzrůstá nepřehlednost.
Také rozšiřování zaměření nakonec vedlo k nutnosti rozdělení, aby byla zachována přehlednost a funkcionalita.
Proto existují specifikace pro určitá odvětví:

\begin{itemize}
	\item \textbf{Enterprise} - síťové a cloudové technologie,
	\item \textbf{ICS} - průmyslové řídící systémy,
	\item \textbf{Mobile} - mobilní zařízení.
\end{itemize}

\paragraph{}
Samozřejmě i jednotlivá odvětví jsou dále rozdělena na domény dle zaměření útočníků.
U Enterprise je rozlišování pro Windows, macOS nebo Linux a také cloud, network a containers.
Mobilní platforma odlišuje Android a iOS\@.
Toto rozdělení dává smysl, stejná technika bude jinak aplikována na různých systémech a díky tomuto rozdělení je možné zamířit se právě na jeden určitý use-case útoku.

\subsubsection{Taktika}
Taktika představuje nejvyšší úroveň abstrakce modelu ATT\&CK\@.
Představuje cíl útočníka, ne však jak se tohoto cíle útočník snaží dosáhnout.
Příkladem můžeme zmínit \textit{Impact(Dopad)}.
Taktika dopadu je zaměřená na způsobení škody v podobě ohrožení integrity nebo dostupnosti dat.
Spadají do této taktiky techniky zaměřené na ničení nebo manipulaci dat.
Rozdělení na taktiky není ojedinělé, naopak na této abstraktní úrovni ATT\&CK těsně připomíná další modely hrozeb zaměřené na postup útoku a útočníka.
Důvod, proč je ATT\&CK preferován před jinými modely je jeho rozsah, provázanost a modelování namapované na proces navigace útočníka v napadaném systému.
Další výhodou je jasná rozlišitelnost taktik.
Jsou označené jedinečným řetězcem ve tvaru \textbf{TA0XXX} kde XXX je nahrazeno číslem jednotlivé taktiky.
Pozor na fakt, že taktika se stejným názvem, například \textit{Initial access(Počáteční přístup)}, bude označena jiným řetězcem v závislosti na odvětví do kterého spadá.
Pro enterprise je to \textbf{TA0001}, pro mobile \textbf{TA0027} a pro ICS se tato taktika značí \textbf{TA0108}.
Nelze tak dojít k záměně a nahlížet na techniky, které nejsou aplikovatelné na dané odvětví.\cite{Mitre_finding_cyber_threats, mitre_attack_framework}

%todo add taktiky z mitre enterprise, ics?

\subsubsection{Technika a sub-technika}
Techniky jsou způsob, kterým útočník dosahuje svého taktického cíle.
To znamená, že v rámci každé taktiky existuje konečný počet popsaných akcí, jejichž provedením útočník dosáhne svého cíle.
Je jasné, že použité techniky se mění v závislosti na znalostech i možnostech útočníka.
Také je nutné brát v potaz prostředí, kde útok probíhá.
Popis taktik je však od těchto specifik oproštěn.
Důvodem je opět nutnost klasifikace, kdy není možné vytvořit téměř nekonečné množství technik pro jednotlivé útoky.
Jsou proto klasifikovány do kategorií technik, kde změny v metodologii nebo jednotlivé kroky jsou přidány pomocí sub-technik.
Tento postup umožňuje popsat nejdůležitější kroky chování útočníka aniž by se obránce ztratil v detailech.
Další výhodou je izolace od neustále vyvíjejícího kodu.
Základní metodika zůstává zachována, avšak jednotlivé implementace se v čase rapidně mění.
Stejně jako taktiky i techniky jsou označeny jedinečným řetězcem, který je ve tvary \textbf{TXXXX} kde XXXX je nahrazeno číslem jednotlivé techniky.
Pokud pro tuto techniku existují sub-techniky, připojuje se na konec řetězce \textbf{.XXX} jde XXX reprezentuje jednotlivou subtechniku.
Za příklad je možno uvézt techniku \textit{Phishing}, která je identifikována řetězcem \textbf{T1566}.
Technika phishing má také 3 sub-techniky, kterými jsou \textit{Spearphishing Attachment \textbf{[T1566.001]}}, \textit{Spearphishing Link \textbf{[T1566.002]}} a \textit{Spearphishing via Service \textbf{[T1566.003]}}


\paragraph{}
Techniky se jen málokdy vyskytují v izolaci, většinou na sebe stejně jako celé taktiky navazují.
Vytvářejí tak sled událostí, které vedou k úspěšnému útoku.
Schopnost identifikovat a propojit navazující druhy událostí je důležitou součástí obrany.
Cílem modelu ATT\&CK je totu schopnost co nejvíce usnadnit.
Data získaná z každého kroku v sekvenci pak nejen vypovídají o tom, zda se jedná o útok ale také obránci dovolují lépe vyhledávat další postup útočníka.
Je také možné, že podle cíle techniky je daná technika zařazena do více taktik zároveň.
Jedním takovýmto příkladem může být technika \textit{input capture(zachycení vstupy)}.
Ta spadá ve stejnou chvíli do taktik \textit{Collection(Sběr dat)} a \textit{Credential access(Pověření k přístupu)}.
Je jasné, že framnework ATT\&CK nevypisuje všechny existující techniky útoku z dané taktiky.
Jeho založení na informacích získaných z opravdových bezpečnostních incidentů a komunity výzkumníků však vede k upřednostnění nejznámějších a nejpoužívanějších technik.
Zdá se to být nejlepší způsob, jak shromažďovat účinné informace o použití technik útočníka a jak je propojovat do užitečného analytického nástroje.\cite{Mitre_finding_cyber_threats, mitre_attack_framework}

\subsubsection{Další přínos}
V neposlední řade ATT\&CK také sleduje původ útoků, využitý software a zařazuje útoky do předpokládaných kampaní.
Tyto záložky jsou na hlavní stránce, specificky se jedná o \textit{Groups, Software} a \textit{Campaigns}.

\paragraph{Groups - skupiny}
Skupiny označují shluky útočníků, kteří operují se stejným cílem a je rozumná domněnka, že jsou mezi sebou v kontaktu.
Analytici sledují skupiny pomocí různých  metodologií a termínů.
Je proto možné, že různé skupiny analytiků nazývají jistou skupinu různými jmény, ATT\&CK se snaží tyto různá pojmenování respektovat.
Proto jsou u skupin zaznamenávány i alternativní názvy.
Také mezi jednotlivými skupinami dochází k částečnému překrytí, a to z důvodu pohyby útočníků a složitosti sledování skupin.
U kriminálních skupin je častý takzvaný \textit{rebranding}, kdy skupina změní své oficiální jméno nebo identifikační rysi za účelem zmatení autorit.
Pro zmatení ja také časté použití identifikačních rysů jiných skupin.
Proto je zaznamenávání a sledování obtížné.
Hlavním přínosem sledování skupin jsou informace o postupech a metodologiích útoku, které daná skupina preferuje.
Tyto informace jsou získány mapováním na veřejně oznámené použití technik a jsou uvedeny původní odkazy.
Je důležité zmínit, že informace nezobrazují všechny možné použití technik skupin, ale pouze podmnožinu, která byla zaznamenána.
I tak je to pro obránce značný přínos.
Při identifikaci metodologie útoku na určitou skupinu je možné předpovídat další kroky z předchozích záznamů a zmírnit tak dopad útoku.\cite{Mitre_groups}


\paragraph{Software}
Software je termín pro vlastní nebo komerční kod a nástroje používané útočníky pro provádění technik nebo taktik.
Stejně jako u skupin můžeme u software najít alternativní názvy právě díky sledování několika skupin.
Překrývání je týmem ATT\&CK a tyto názvy jsou označeny jako \textit{Associated Software}.
Informace o software pochází z veřejně oznámených technik a často bývají mapované na skupiny, které tento software využívají.
Je tedy opět nutné zmínit, že použití určitého software neznamená, že se jedná o související skupinu.
Pouze to znamená, že daná skupina tento software v minulosti použila.\cite{Mitre_software}


\paragraph{Campaigns - kampaně}
Pojem kampaň označuje rozsáhlou intruzivní aktivitu, která proběhla během určitého období s obecnými cíli a úkoly a byla značnou kybernetickou hrozbou.
Jedná se například o dlouhodobou průmyslovou špionáž nebo špionáž proti vládě určité země.
není neobvyklé, že do složité kybernetické operace je zapojeno několik přidružených skupin, přičemž každá hraje jedinečnou roli.
Například může být jedna skupina zodpovědná za počáteční průnik do systému a jiná za exfiltraci dat
Kampaň je označována názvem uvedeným ve veřejných zprávách nebo jedinečným identifikátorem ATT\&CK pokud kampaň ještě nebyla pojmenována.
Zaznamenání a sledování kampaní je obtížné.
Různé organizace zaznamenávající útok patřící do kampaně ho mohou oznámit pod různými jmény, a stejný útok navíc může být popsán z různých úhlů pohledu.
Tým ATT\&CK pro jednotlivé kampaně činí nejlepší úsilí o sledování překrývajících se jmen, která jsou označena jako "Přidružené kampaně".
Toto sledování je výhodné, protože poskytuje informace nejen o existenci kampaně, ale také informace o technikách, které jsou využívány.
Pro referenci jsou tak kampaně mapovány na veřejně uváděné techniky a jsou v nich zahrnuty původní reference.\cite{Mitre_campaigns}

\subsubsection{Nástroje}
~\cntcapfigure{cso_mitre_open-source_attck_test_tools}{\linewidth}{Porovnání vybraných nástrojů}{\cite{csoonline_4_testing_frameworks}}\label{fig:OpenSourceTestTools}
% TODO maybe make my own table?

Než je zvolen finální framework pro testování, je nutné si nejdříve připravit co možná nejvíce informací týkajících se testování.
Je jasné, že každý z vybraných frameworků má své silné a slabé stránky.
Jak je vidět v tabulce ~\ref{fig:OpenSourceTestTools}, jedním z ovlivňujících parametrů je například testovaný operační systém.
Kde s atomic red je možné testovat na všech třech hlavních platformách, \ac{RTA} je přizpůsobené pouze pro Windows.
Specializace sama o sobě není špatná vlastnost, naopak většinou bývá přínosem.
Pokud však potřebujeme testovat Linux, je jasné že nemůžeme použít nástroj který toho není schopen.
Žádný z frameworků také nepokrývá celou MITRE ATT\&CK matice.
Proto lze předpokládat použití více nástrojů právě podle okruhu testování, na který je testování zaměřeno.

\subsubsection{Red Canary Atomic Red}
Ze všech uváděných příkladů je \textit{Atomic Red} nejvíce přívětivý rychlému začátku .
Pro jeho funkčnost není potřeba nic stahovat ani nastavovat, lze ho jednoduše použít z příkazové řádky\cite{redcanary_atomic_red}.
Další odlišnost je jeho ovládání.
Většina ostatních frameworků je založena na jazyku Python, popřípadě jiném skriptovacím jazyce.
Atomic red je ale spíše sada instrukcí odpovídajících jednotlivým taktikám a technikám Mitre ATT\&CK\@.
Atomic Red tedy není testovací software ale instrukce, jak testovat jednotlivé ATT\&CK taktiky s nástroji, které jsou již pravděpodobně nainstalované v počítači.
Je tedy vhodný pro naučení testovacích postupů ne však pro automatizované testování.
Uživatel při testování vždy musí provézt kroky útoku a sledovat výsledek.
To neznamená, že není možné simulovat složité útoky využívající hned několik technik.
Atomic red dovoluje techniky za sebou řetězit a vytvářet složité sequence útoků.
Výborným příkladem přímo od tvůrců je kod zvaný \textit{Dragons tail}\cite{redcanary_dragons_tail}.
Tento příklad se snaží napodobit chování skutečných útočníků.
Samozřejmě napodobení není dokonalé, ale dobře ukazuje možnosti Atomic Red pro složitější testování.\cite{csoonline_4_testing_frameworks}

\paragraph{Nevýhody}
Hlavní nedokonalostí tohoto přístupu je nutnost vytvářet nebo najít vlastní pomůcky.
Atomic Red je nástroj pro obránce a nechce útočníkům všechny potřebné prostředky pro provedení útoku.
To znamená že pomocné nástroje pro útok si musí obránce připravit sám.
Příkladem může být malware vložený v Microsoft Word dokumentu.
Atomic Red se snaží ukázat, jak malware operuje a jek ho útočníci využívají bez toho, aby ho útočníkům poskytl.
Dalším problémem je generace výstupních reportů.
Na rozdíl od dalších přístupů Atomic Red nevytváří automaticky stručný report o technikách které byly v útoku zkoumány.
Tato funkcionalita je bohužel až v pokročilejším placeném software který Red Canary nabízí.\cite{csoonline_4_testing_frameworks}

%todo add more after using it?
% TODO WRITER: text
%todo co to je, porovnání

\subsubsection{Endgame RTA}
\ac{RTA} v tomto kontextu znamená \firstac{RTA}.
Z představených nástrojů je to stále jednodušší, i když mocný nástroj.
Jediný povinný požadavek je Python verze 2.7, kde pro instalaci stačí rozbalit z gitu stažený repozitář.
Pro plnou funkcionalitu je však doporučeno navíc do podadresáře \textit{bin} vložit Sysinternals Suite a MsXsl.
Stále je to však jednoduchá instalace a celý návod se nachází přímo v git repozitáři.\cite{endgame_rta}

\paragraph{Schopnosti}
\ac{RTA} je složení téměř 50 různých skriptů simulujících útoky.
Bohužel jména těchto skriptů neodpovídají přesně Mitre ATT\&CK, i když jsou na ni založené.
Také zde neexistuje hyperlink navigace jako například u Atomic Red.
Práce se skripty proto vyžaduje více úsilí a určitou známost matice ATT\&CK\@.
Naštěstí jsou přesná označení technik, jako T1107 dohledatelné přímo ve skriptu.
Pro složitější testování je možné spouště všechny skripty najednou nebo vybrat některé pomocí podmíněného příkazu \textit{IF} přímo v kodu.
Pro základní automatizaci stačí, pro hlubší testovaní je tento nástroj však omezen.\cite{csoonline_4_testing_frameworks}

%todo probably better to use never EQL analytics library? https://eqllib.readthedocs.io/en/latest/
% TODO WRITER: text


\subsubsection{Mitre Caldera}
Caldera je od stejných vývojářů jako ATT\&CK, tedy Mitre.
Lze tak očekávat dobrou integraci a kompatibilitu mezi těmito nástroji.
Na rozdíl od dříve představených nástrojů je Mitre Caldera komparativně náročnější nástroj, což se projevili i v instalaci.
Ani tak však není instalace náročná, a to hlavně díky přípravě tvůrců a dobré dokumentaci\cite{mitre_caldera_docs}.
Framework caldera lze rozdělit na 2 hlavní části.
První je \textit{Jádro systému}.
Tento framework je dostupný v github repozitáři.
Zahrnuje asynchronní server pro řízení příkazů a kontrolu (command-and-control, C2) s REST API a webovým rozhraním.
druhou složkou jsou \textit{pluginy}.
Jsou to rozšíření pro možnosti jádra frameworku poskytující dodatečnou funkcionalitu.
Většinou se jedná o samostatné repozitáře.
Jako příklady je dobré uvézt agenty, hlášení, nebo  techniky z Atomic red team.\cite{mitre_caldera}

\paragraph{Systémové požadavky}
Autoři se snaží podporovat široké spektrum cílových systémů.
Proto existují 2 typy požadavků, povinné a doporučené.
Povinné určují, kde může být vůbec Caldera připravena.
Jádro frameworku může být nainstalováno na operačních systémech Linux nebo MacOS.
Pro fungování je vyžadován Python verze 3.7, 3.8, nebo 3.9 samozřejmě s rozhraním pip3.
Pip3 je vyžadován pro instalaci podpůrných balíčků uvedených v souboru \textit{requirements} přímo v repozitáři Caldery.
Ke správnému ovládání je vyžadován moderní prohlížeč.
Doporučený je Google Chrome.
Doporučené požadavky se týkají vývoje a správného fungování Caldery.
Je doporučený hardware s 8GB+ RAM a 2+ CPU\@.
Dále je pro správnou kompilaci a funkčnost Golang agentů doporučen GoLang 1.17+.
Za zmínku stojí také možnost instalovat a pracovat s Calderou v Docker kontejneru.\cite{mitre_caldera_docs}

\paragraph{Možnosti}
Caldera je oproti dříve představeným nástrojům složitý framework.
To je však více nežli dostatečně vyváženo dobrou dokumentací a hlavně možnostmi, které zkušenému uživateli Caldera nabízí.
Caldera dovoluje nastavení autonomních red-team nebo blue-team operací, a to plnohodnotně včetně závěrečného výpisu a exportování logů.
Existuje také možnost manuálních red-team operací, hlavně v případech nahrazování a testování pomocí vlastních nástrojů.\cite{mitre_caldera_docs}
%todo add more after using it?



\section{Praktická část}

% TODO WRITER: text

\subsection{Popis testovacího prostředí}

kali, kali purple, instalace caldebry, využití vmvare, proč, co a jak
zameření na přenositelnost a rychlou instalaci

\subsection{Popis testované infrastruktory}

založeno na nemocniční prostředí, poradit se s honzíkem co můžem nemužem napsat, co tam je, co tam není atd atd

říci jak je to implementováno, takže lab -> jaké stroje, jak to je propojené atd

\subsection{Testované scénáře}

přidat jsem identifikované vulnarability z navigátoru, u každého scénáře pak z navigáto přesně to, co je testované

\subsubsection{Scénář 1 - doplnit název}

\subsubsection{Scénář 2 - doplnit název}

\subsubsection{Scénář 3 - doplnit název}

\subsection{Celkové zhodnocení výsledků}

\section{Závěry a doporučení}

% TODO WRITER: text

Kritická diskuze nad výsledky, ke kterým autor dospěl (soulad výsledků  literaturou či předpoklady;
výsledky a okolnosti, které zvláště ovlivnily předkládanou práci atd.).
Je vhodné naznačit i případné další
(popř. alternativní) možnosti zkoumání dané problematiky a otevřené problémy pro další studium.

% todo del - this is example and testing section

Vlastní řešení dokládá student zpravidla v několika kapitolách.
Podle charakteru práce musí student uvážit, zda informace
netextové povahy (data, tabulky, obrázky atd.) bude uvádět přímo v textu, nebo je zařadí až za celou práci ve formě příloh, či bude kombinovat oba způsoby.
Více podrobností viz Metodické pokyny pro vypracování bakalářských a diplomových prací (zveřejňované formou výnosů děkana)
a v kurzu MES – Metodologický seminář.

	\subsection{Podkapitola}

	Text podkapitoly.

	Následuje ukázka nějakého seznamu:
	\begin{itemize}
		\item fotografie/avatar,
		\item kategorizační štítky,
		\item popisek.
	\end{itemize}

		\subsubsection{Podpodkapitola}

		Lorem ipsum dolor sit amet, consectetur adipiscing elit. Phasellus sit amet ornare diam, id consequat diam.

			\nlparagraph{Paragraf}

			\noindent Ukázka prvního odstavce v paragrafu. Lorem ipsum dolor sit amet, consectetur adipiscing elit. Phasellus sit amet ornare diam, id consequat diam.



			Použití obrázku je jednoduché, zadáme název souboru, šířku, popisek a zdroj:\cntcapfigure{rovnovaha_paka}{8cm}{Páka rovnováhy vzhledu.}
			V případně, že autor je i autor obrázku:

			\cntcapfigure{uhk}{\linewidth}{Toto je UHK.}{[autor]}

			Blok kódu může vypadat takto:

			\begin{codeblock}
				\begin{verbatim}
@Mapper
public interface AccountDao {
  @Select({
    "select *",
    "from " + Account.TABLE_NAME,
    "where email = #{email}"
  })
  Optional<Account> findByEmail(String email);
}
				\end{verbatim}
				\captionsource{Ukázka bloku kódu.}{[autor]}
			\end{codeblock}

			Tabulka zas může vypadat takto:

			\begin{table}[hbt!]
				\captionsource{Ukázková tabulka.}{[autor]}
				\centering
				\begin{tabular}{| l | r | r | r | }
					\hline
					&        psnr &      ssim &      doba  \\
					model &       (db)    &           & gen. (s) \\
					\hline
					bik. int. & 28.3155 & 0.8566 & 0.0322 \\
					nn1000    & 30.1461 & 0.9043 & 0.8109 \\
					nn1001    & 30.0324 & 0.9023 & 0.7486 \\
					nn1002    & \textbf{30.1886} & \textbf{0.9046} & 1.1731 \\
					nn1003    & 30.0390 & 0.9030 & 1.1320 \\
					nn1004    & 24.9772 & 0.7172 & 4.4367 \\
					nn1005    & 26.1629 & 0.8004 & 4.0475 \\
					nn1006    & 27.9129 & 0.8438 & 4.0683 \\
					nn1007    & 27.5834 & 0.8360 & 4.2082 \\
					\hline
				\end{tabular}
			\end{table}

			\newpage

			Další text

\section{Závěry a doporučení}

% TODO WRITER: text

Kritická diskuze nad výsledky, ke kterým autor dospěl (soulad výsledků  literaturou či předpoklady;
výsledky a okolnosti, které zvláště ovlivnily předkládanou práci atd.). Je vhodné naznačit i případné další
(popř. alternativní) možnosti zkoumání dané problematiky a otevřené problémy pro další studium.


  \iftwoside
    \cleardoublepage
  \fi

  % Pouzita literatura / Bibliography
  % ----------------------------------------------

\ifslovak
  \makeatletter
  \def\@openbib@code{\addcontentsline{toc}{section}{Literatúra}}
\else
  \ifczech
    \makeatletter
    \def\@openbib@code{\addcontentsline{toc}{section}{Seznam použitých zdrojů}}
  \else
    \makeatletter
    \def\@openbib@code{\addcontentsline{toc}{section}{Bibliography}}
  \fi
\fi
  \makeatother
  \begin{flushleft}
  \label{references}

  \bibliographystyle{bibstyle}
  \bibliography{literature}
  \end{flushleft}

  \newpage

  % vynechani stranky v oboustrannem rezimu
  % Skip the page in the two-sided mode
  \iftwoside
    \cleardoublepage
  \fi

  \ifslovak
   \addcontentsline{toc}{section}{Zoznam zkratek}
  \else
    \ifczech
      \addcontentsline{toc}{section}{Seznam zkratek}
    \else
      \addcontentsline{toc}{section}{List of Abbreviations}
    \fi
  \fi
  \startcontents[section]
  \label{abb}
  \setlength{\parskip}{0pt}
  % seznam příloh / list of appendices
  % \printcontents[section]{l}{0}{\setcounter{tocdepth}{2}}

  \printacronyms[name=Seznam zkratek,heading=section*]

  \newpage
  % vynechani stranky v oboustrannem rezimu
  \iftwoside
    \cleardoublepage
  \fi

  % Seznam obrazku a tabulek (pokud prace obsahuje velke mnozstvi obrazku, tak se to hodi)
  % List of figures and list of tables (if the thesis contains a lot of pictures, it is good)

  \addcontentsline{toc}{section}{Seznam obrázků}
  \ifczech
  \renewcommand\listfigurename{Seznam obrázků}
  \fi
  \ifslovak
  \renewcommand\listfigurename{Zoznam obrázkov}
  \fi
  {\hypersetup{hidelinks}\listoffigures}
  \thispagestyle{empty}

  \newpage

  \addcontentsline{toc}{section}{Seznam ukázek kódů}
  \listofcodeblocks

  \newpage

  %\includepdf[pages=-, offset=20 -50, width=20cm]{Images/assignment.png}% this is original
  \includepdf[pages=-, offset=20 -40, width=18cm]{Images/assignment.png}

\end{document}
